\section{Introduction}

\subsection{Background of the Study}

Student academic performance is influenced by a combination of academic, personal, and environmental factors. Understanding these factors can help educators identify students who may require additional support and develop interventions that improve learning outcomes.

Statistical modelling provides a useful framework for studying these relationships. By analysing measurable characteristics such as study time, attendance, previous academic performance, and access to learning resources, it is possible to investigate how different factors are associated with students' academic outcomes.

\subsection{Problem Statement}

Educational institutions often collect information about students, but the relationships between different characteristics of students and academic performance may not be immediately apparent. A statistical model can help identify important predictors of performance and quantify their relationships with academic outcomes.

\subsection{Aim of the Study}

The purpose of this study is to develop a statistical model to examine factors associated with student academic performance.

\subsection{Research Objectives}

The specific objectives are to:

\begin{enumerate}
    \item identify factors that are associated with academic performance;
    \item examine the relationship between study behavior and academic outcomes;
    \item develop a statistical model for predicting academic performance; and
    \item evaluate the performance of the proposed model.
\end{enumerate}

\subsection{Research Questions}

The study addresses the following questions:

\begin{enumerate}
    \item What factors are significantly associated with academic performance of the student?
    \item How does study time relate to academic performance?
    \item Can academic performance be predicted using measurable student characteristics?
\end{enumerate}


 